
\begin{proposition}[The unrestricted normalization target under the integral CDP hypotheses]\label{cdp:st:C4e}
Let $f:Y\to\PP^1$ satisfy the integral CDP hypotheses and $c_3(Y)>0$. Then a fixed irreducible nonnormal reduced support $S$ has normalization $\eta:\widetilde S\to S$ satisfying
\[
 H^2\bigl(\widetilde S,\omega_{\widetilde S}\otimes
           \eta^*(K_Y^{-1}\otimes L)\bigr)\ne0
 \qquad\hbox{for every }L\in\Pic(Y).                  \tag{C.21}\label{cdp:eq:C.21}
\]
The assertion concerns the reduced support; no multiplicity-one conclusion is included.
\end{proposition}
\begin{proof}
\ref{cdp:st:C.3.1} and the exclusions \ref{cdp:st:C.3.7}, \ref{cdp:st:C.3.9} and \ref{cdp:st:C.3.10} show
that the general fibre is a two-torus, a Hopf surface, or
an Inoue surface. In the last two cases it has
$h^1(\OO)=1$ and $h^2(\OO)=0$.
For torus general fibre, \ref{cdp:st:C4c} gives a fixed full fibre with a nonzero ambient cotangent section for every ambient twist. In the second case, \ref{cdp:st:C4d} gives persistence along $\C\beta$, and \ref{cdp:st:A.5.13} identifies this with $\Pic(Y)$. From this point the proof uses precisely that persistence property and does not use $c_3(Y)>0$ again. \ref{cdp:st:A.3.5} makes its support irreducible, so it is $mS$ for a single positive integer $m$. \ref{cdp:st:A.3.4} excludes a smooth support, and \ref{cdp:st:A.5.2} then excludes a normal support. Its filtration in \ref{cdp:st:A.3.1} has finitely many graded factors
\[
 \Omega_Y^1|_S\otimes M|_S\otimes\OO_Y(-jS)|_S,
 \qquad 0\le j<m.
\]
Nonvanishing on the filtered sheaf requires nonvanishing on at least one factor. Each factor's nonvanishing locus is closed analytic in the connected vector group $\Pic(Y)$. A finite union of proper such loci cannot cover that group. Hence one factor has nonzero sections for every $M$. Translating by $\OO_Y(-jS)$ shows
$H^0(S,\Omega_Y^1|_S\otimes M|_S)\ne0$ for every $M$.
\ref{cdp:st:A.4.9} kills the torsion-free quotient on a dense open set of twists. The exact sequence \ref{cdp:st:A.1.1} and proper semicontinuity then show that the inverse-Jacobian kernel itself has sections for every twist.

We apply the precise persistent-kernel reduction \ref{cdp:st:A.5.4}. It gives a fixed component of the normalization on whose resolution every twist of one fixed rank-one sheaf has a section. The Gauduchon degree argument there proves that the pullback of every ambient line has degree zero for every Gauduchon metric. Its conductor comparison excludes nonnegative Kodaira dimension, and its class VII argument makes every ambient line trivial on that normalization component. Rational components have the same triviality because their Picard group has no connected part. Either case gives \eqref{cdp:eq:C.21} immediately by Serre duality and descent from resolution. It remains to treat an irrational ruled resolution $p:Z\to B$.

Let $h:Z\to S\subset Y$. \ref{cdp:st:A.5.3} supplies an effective conductor divisor $D_Z$ with
\[
 K_Z+D_Z=h^*(K_Y+S),\qquad D_Z\cdot F=2              \tag{C.22}\label{cdp:eq:C.22}
\]
for a general ruling fibre $F$. Its horizontal part is therefore either reduced of total ruling degree two, or $2C$ for a section $C$. We record the comparison through each ruling blowdown $\sigma:Z\to Z'$ used below. Put $D'=\sigma_*D_Z$, let $a$ be the coefficient of its exceptional curve $E$ in $D_Z$, and put $d=\operatorname{mult}_{\sigma(E)}D'$. Then
\[
 K_Z+D_Z=\sigma^*(K_{Z'}+D')+(1+a-d)E.
 \tag{C.22b}\label{cdp:eq:C.22b}
\]
Intersection with $E$, using $c_1(K_Z+D_Z)=0$ and $E^2=-1$, gives $a=d-1$. Thus the log-canonical divisor is the exact pullback of $K_{Z'}+D'$, and the latter again has zero rational Chern class. No factorization of $h$ through the blowdown is asserted or needed.

We first treat the reduced horizontal part $H$. If a component is exceptional for the normalization resolution, every ambient line restricts trivially to it. Its normalization maps finitely onto $B$, and the norm composed with pullback on $\Pic^0(B)$ is multiplication by the covering degree. Hence this restriction has finite kernel; the connected ambient Picard image is trivial, giving \eqref{cdp:eq:C.21}. We may therefore assume that the horizontal components map to conductor curves. Adjunction on each of them gives
\[
 2p_a(H_i)-2+(D_Z-H_i)\cdot H_i=0.
\]
The second term is nonnegative. Since $H_i$ dominates a positive-genus base, its arithmetic genus is at least one. Thus the base has genus one, each $H_i$ is smooth elliptic, and the horizontal curves are disjoint from each other and from the vertical part. Contract ruling $(-1)$-curves to a minimal bundle. A contracted curve outside $D_Z$ has $D_Z\cdot E=1$, so cannot identify two horizontal points or create a horizontal node. A contracted curve in the vertical part is disjoint from $H$. Downstairs the horizontal curves remain smooth and disjoint, and any nonzero vertical divisor on a minimal ruled bundle meets them. Consequently the downstairs vertical part is zero. Reversing the contractions, a blowup away from the horizontal conductor would give coefficient $-1$ on its exceptional curve in \eqref{cdp:eq:C.22}, which persists under further strict transform and contradicts effectiveness. A blowup at a smooth point of a coefficient-one horizontal curve gives exceptional coefficient zero. Hence $D_Z=H$ has no vertical component. Every codimension-one singularity of $S$ is then an ordinary node, and \ref{cdp:st:A.2.3} identifies the inverse-Jacobian kernel with the normalization kernel globally, including its codimension-two points. A section of a topologically trivial line on the Kahler resolution has an effective zero divisor of zero Kahler volume, so is nowhere zero. Persistence for all twists therefore makes all ambient pullbacks trivial and proves \eqref{cdp:eq:C.21}.

We now have $D_Z=2C+V$, where $C$ is a section and $V$ is vertical. If $C$ is exceptional for $Z\to\widetilde S$, every ambient line is trivial on $C$. Restriction $\Pic^0(Z)\to\Pic^0(C)$ is an isomorphism, because $C$ is a ruling section. This again makes all ambient pullbacks trivial. Otherwise let $D\subset S$ be its image and let $D^\nu$ normalize that curve. When $D^\nu$ is rational, every ambient line pulls back trivially to $D^\nu$, hence to $C$, giving the same conclusion.

Suppose instead that $g(D^\nu)\ge1$. A vertical rational curve cannot map finitely onto $D^\nu$, so all the normalization branches over the general point of $D$ belong to $C$. Their conductor coefficient is two. We derive the possible transverse reduced plane curves rather than assuming a list of singularities. Write its equation as $f=\prod_i f_i$ in the smooth-plane local ring and let $R_i$ and $B_i$ be a branch ring and its normalization. The elements of $R=\OO/(f)$ supported only on branch $i$ form
\[
 \left(\prod_{j\ne i}f_j\right)\bigg|_{f_i}\,R_i.
\]
This equality follows from
$\bigcap_{j\ne i}(f_j)=(\prod_{j\ne i}f_j)$ in that unique
factorization domain. Multiplying by arbitrary elements of
$B_i$ therefore shows that the conductor exponent on branch $i$
is its intrinsic exponent plus $\sum_{j\ne i}I(f_i,f_j)$.
A singular branch has intrinsic exponent at least two:
exponent one would imply $\C+t\C\{t\}=\C\{t\}\subset R_i$.
Exponent two forces $R_i=\C+t^2\C\{t\}=\C\{t^2,t^3\}$,
since an element of order one would again make the ring smooth.
Any additional branch has intersection multiplicity at least
two with a singular branch and is therefore excluded. Thus a
singular branch gives exactly the cusp. When all branches are
smooth, their pairwise positive intersection multiplicities sum
to two on each branch. This gives two branches of contact two
(the tacnode), or three pairwise transverse branches (the
ordinary triple point), and exhausts the possibilities.
The only remaining alternative is a nonconstant pullback
$\Pic(Y)\to\Pic(Z)$; a constant pullback is the zero
homomorphism and already gives \eqref{cdp:eq:C.21}. In that alternative,
the coherent conductor argument \ref{cdp:st:A.5.12} gives $p_a(D)\le1$
for torus general fibre, and \ref{cdp:st:A.5.13} gives the same bound
in the second case.
Since $g(D^\nu)\ge1$, it follows that $D$ itself is smooth
elliptic. This conclusion permits nonreduced conductor image
schemes and imposes no assumption on the original singular
curve. We can therefore apply \ref{cdp:st:C2d} directly. It makes the
transverse type constant, gives the etale conductor cover
$C\to D$, and proves $C^2=0$. In particular $C$ and the
ruling base are elliptic.

Adjunction in \eqref{cdp:eq:C.22} gives $V\cdot C=0$. No ruling $(-1)$-curve meets $C$: for one outside $V$, its intersection with $K_Z+2C+V$ would be at least $-1+2=1$; for one in $V$, the equality $V\cdot C=0$ forbids such an intersection. Thus every contraction to the minimal ruled bundle avoids $C$ and preserves its self-intersection zero. On that bundle $K+2C$ is numerically trivial, so the effective vertical part is numerically trivial and hence zero. A first reverse blowup, necessarily away from $C$, would again produce negative conductor coefficient. There are therefore no reverse blowups: $Z$ is minimal ruled and $D_Z=2C$. There is also no exceptional curve for the normalization resolution: a negative curve $E\ne C$ would satisfy $2p_a(E)-2=E^2-2C\cdot E<0$, hence would be rational and contained in a ruling fibre; a minimal ruled bundle has only self-intersection-zero fibres. Thus the normalization itself is smooth here.

By the explicit inverse-Jacobian computations in \ref{cdp:st:C2d}, its pullback allows at most a simple pole along $C$. There are no other divisorial poles. Hartogs extension for the line bundle on the smooth surface gives an injection of the kernel into $h^*\OO_Y(-S)\otimes\OO_Z(C)$. The exact sequence of the section $C$ and $R^1p_*\OO_Z=0$ give
\[
 0\longrightarrow\OO_B\longrightarrow p_*\OO_Z(C)
   \longrightarrow N_{C/Z}\longrightarrow0,
 \qquad \deg N_{C/Z}=0.                              \tag{C.22a}\label{cdp:eq:C.22a}
\]
For a degree-zero line $P$ on the elliptic curve, both end terms after tensoring by $P$ have zero sections unless $P=\OO_B$ or $P=N_{C/Z}^{-1}$. Thus $H^0(Z,\OO_Z(C)\otimes p^*P)$ is supported at at most two points of $\Pic^0(B)$. The ambient pullback image is a connected complex subgroup of this Picard group. Persistence for every ambient twist forces that subgroup to be a point, and it contains the identity. Therefore every ambient line pulls back trivially to $Z$ in this last case as well.

The equality $\pi_*\OO_Z=\OO_{\widetilde S}$ for the normal resolution descends triviality of a pulled-back line: projection formula identifies its space of sections with that of its trivial pullback, and the same statement for its inverse gives mutually inverse trivializations. Consequently one normalization component has trivial $\eta^*(K_Y^{-1}\otimes L)$ for every $L$. Serre duality on that compact normal Cohen--Macaulay surface identifies its top dualizing cohomology with $H^0(\OO)^\vee=\C$. It is a direct summand of the target in \eqref{cdp:eq:C.21}, proving the claim.
\end{proof}

The local example \ref{cdp:st:A.2.5} has actual vertical inverse-Jacobian
poles, and \ref{cdp:st:A.2.6} has a critical curve that is not a local
complete intersection. Their local rings remain part of the
audit. The global conductor-duality arguments \ref{cdp:st:A.5.12}--\ref{cdp:st:A.5.13} prove
that a positive-genus conductor image in the persistent ruled
case is smooth, and the exact contractions above then remove
the vertical conductor. Thus neither local phenomenon leaves
an unproved step in the stated theorem. The assertion
is about the normalization of a reduced support and does not
claim that its full fibre has multiplicity one.


